\section*{5. Conclusion}

The optimization model provides Universal Airways with a well-structured plan for fleet expansion and route allocation. Based on the analysis, it is recommended that Universal Airways purchase six Med-848 planes and four Smal-286 planes. The Med-848 planes, which offer higher revenue per flight, are best utilized on the high-demand, long-haul routes to Metropolis and Cityburg. The allocation of eight daily flights to each of these cities ensures maximum revenue generation by fully utilizing these aircraft for the most profitable destinations. This approach allows the airline to capitalize on the strong demand for these routes while maintaining operational efficiency.

On the other hand, the Smal-286 planes, being less expensive and more suitable for shorter routes, are ideal for servicing destinations with lower passenger demand, such as Smalltown and Townville. The analysis recommends allocating five daily flights to Cityburg, three to Smalltown, and four to Townville. This allocation meets the operational requirement to maintain service to these regions, even if they are less profitable. Additionally, the Smal-286 planes' ability to service shorter flights efficiently ensures that the airline can provide reliable service while keeping costs under control.

In summary, the proposed fleet composition and flight schedule are designed to maximize daily profit while adhering to budgetary and operational constraints. By purchasing six Med-848 planes and four Smal-286 planes, Universal Airways can optimize its operations to achieve a maximum daily profit of \$160,400. This solution not only meets the company's expansion goals but also provides a framework for sustained profitability as the airline expands its service to new destinations.

While this model offers an optimal solution given the current constraints, future work could extend the analysis by incorporating additional factors. For example, demand uncertainty, which can vary due to seasonal fluctuations, economic conditions, or unforeseen events, could be integrated into the model to provide a more dynamic and flexible flight schedule. Additionally, considering the volatility of fuel prices, which form a significant part of airline operating costs, could further refine profit projections and offer a more robust strategy for cost management. Moreover, with increasing global attention on environmental sustainability, future models could consider regulatory factors such as carbon emissions or fuel efficiency standards. This would allow Universal Airways to not only optimize for profit but also align with long-term environmental goals.

By considering these additional factors in future optimizations, Universal Airways can ensure that its operations remain both profitable and adaptable in an evolving market. This holistic approach will support the company’s strategic objectives of expanding its fleet, maintaining service to new regions, and maximizing profitability.
